\section{Evaluation}
\label{sec:eval}

\subsection{Experimental Setup}

\begin{figure}
\begin{tabular}{ | m{5em} | m{17em} | } 
  \hline
  Frontend & 2-way superscalar, 256 entries BTB, 8 entries RAS, tournament predictor \\
  \hline
  Reorder Buffer &  64 entries, 2-way insert/commit \\
  \hline
  Reservation Stations & 2 stations for ALU (16 entries), 1 for FPU (16 entry), 1 for memory (16 entries) \\
  \hline
  L1 (I/D) &  32KB, 8 ways, 8 outstanding requests \\
  \hline
  Load-Store Unit & 24 LdQ entries, 14 StQ entries, 4 SB entries\\
  \hline
  LLC (L2) & 1MB, 10 cycles latency, 16 ways, max 16 outstanding requests \\
  \hline
  Memory &  120 cycles latency (24 outstanding requests) \\
  \hline
\end{tabular}
    \caption{Citadel's microarchitectural configuration.}
    \label{tab:mi6param}
\end{figure}
Citadel's processor is a modified out-of-order 2-core based on Riscy-OOO~\cite{zhang2018composable} and MI6~\cite{bourgeat2019mi6}.
The parameters used for the cores and the memory subsystem are given in Figure~\ref{tab:mi6param}.
Citadel runs on an AWS EC2 F1 instance, equipped with a Virtex UltraScale+ FPGA. 
On this FPGA board, Citadel can be clocked at $30$ MHz, which is similar to the baseline design.
We use the RISC-V tool chain with GCC 12.2.0.
We compile all software with the -O3 flag. 
We run Linux Kernel 6.2.0.
Unless mentioned differently, all performance numbers are collected by running the benchmark on the FPGA using performance counters and an average of ten runs. 

\subsection{Implementation Data and TCB Breakdown}
\begin{table}
    \centering
    \begin{tabular}{|p{2.35cm}|r|r|}
        \hline
         Component & LOC & Size\\
         \hline\hline
         Processor \cite{bourgeat2019mi6} (+ modification) & $\sim60$ KLOC (+ 250 ) & N/A \\
         \hline
         Bootloader & 5236 & 55KB\\
         \hline
         Security Monitor & 9287 & 97KB\\
         \hline
         \rowcolor{LightRedbis}
         Mini-SM & 2301 & 5.6KB\\
         \hline
         \rowcolor{LightRed}
         Linux & $\sim20$ MLOC & 157MB\\
         \hline
         \rowcolor{LightRed}
         SM Kernel Mod. & 507 & 505KB\\
         \hline
         \hline
         TCB Software & 14523 & 152KB \\
         \hline
    \end{tabular}
    \vspace{0.40cm}
    \caption{Software Component Size and TCB Breakdown. Mini-SM is a subset of the SM code. 
    Lines in red correspond to elements excluded from the TCB.}
    \label{tab:tcbloc}
\end{table}

Number of lines of code (LOC) and binary size of Citadel's components are presented in Table \ref{tab:tcbloc}.
Due to their limited functions, the SM and the Bootloader can be kept small and the total software Trusted Code Base (TCB) is less than 15K LOC.
This helps to minimize potential bugs and reduce the software attack surface.
It can also help to make future formal verification efforts tractable.
Final TCB size needs to include the enclave's binary, which depends on the application. 

\subsection{Evaluating the End-to-End Platform}

\textbf{MicroPython Runtime}
\label{sec:evalmicropython}
We do a minimal port of MicroPython \cite{microphython} inside one of our enclaves.
We start with the bare-metal ARM implementation.
We modify the configuration files and the linker script so our binary can run in the \emph{virtual memory layout} provided by our enclave.
We only need to implement a few functions by hand.
For instance, we redirect the console to print to shared memory and need to write a bit of assembly for the \texttt{setjmp} and \texttt{longjmp} functions used for exception handling.
We also add a cryptographic library, so the enclave can sign the console output using its own private keys recovered through the remote attestation mechanism.
A hash of the standard output is extended on each character print and signed on the enclave exit.
We cross-compile MicroPython using the RISC-V tool-chain, build without standard library, and generate position independent code.
In total, we need to manually modify 128 LOC for the port and add the 5K LOC cryptographic library. 
The final compiled binary measures 180KB.
With the enclave's code added, the TCB is now 65,174 LOC, most of it (70\%) being MicroPython.

\begin{figure}[t]
\centering
\includegraphics[width=\columnwidth]{figures/runtime_breakdown.png}
\vspace{-20pt}
\caption{MicroPython Enclave Execution Breakdown.}
\vspace{-8pt}

\label{fig:micro-pythonenclave}
\end{figure}

We use a simple primality test as a benchmark where a number is found to be prime if none of the numbers smaller than it can divide it. 
We test the primality of an arbitrary small prime number: $19,531$.
The tests take an average of $4.1\times10^7$ cycles to execute, with 51\% of the time spent within the Python program.
This means the enclave and attestation mechanism add 94\% overhead for this small computation.
We show a visual breakdown of the running time in Figure \ref{fig:micro-pythonenclave}.
We note that most of this overhead comes from hashing  enclave code (25\% of the total running time) but also generating the enclave keys and signing the certificate (10\% of the running time). These two numbers are especially high since this code is run non-speculatively inside the SM.
On another hand, hashing and signing the output inside of enclave takes a negligible amount of time (less than 1\% of the total running time).
Because this cost is fixed, it could definitely be amortized for bigger applications.

\noindent\textbf{Securing a Cryptographic Library}
\label{sec:evalcrypto}
We write a wrapper around an ED25519 RISC-V library \cite{ed25519github}.
The wrapper makes it possible to keep the keys used by the library inside of the enclave and for the library functions to be accessed in a Remote Procedure Call style (RPC).
We add a queue to shared memory, so untrusted applications can send requests to the library to generate keys (that will stay in enclave memory), or to sign messages using previously generated keys.
In total, we need to implement 334 LOC of C, which is to be compared to the 4,514 LOC of the crypto library (7\% overhead).
With the enclave's code added, the total TCB is now 19,231 LOC.
The final crypto-enclave binary is 71KB.

We compare performance between three scenarios.
In each of these scenarios, we measure the number of cycles required to sign randomly generated messages.
We split \texttt{1MB} of randomly generated data into $\approx 3000$ packets of various lengths following the arbitrary Internet Mix distribution \cite{imix}.
In the first setup, which corresponds to our insecure baseline, we do not use any enclave as we directly link the cryptographic library to our user-level application.
In the second scenario, called ``multi-core'', the enclave constantly runs on a dedicated core and programs running on other cores can forward requests through the shared memory mechanism.
In the third scenario, termed ``re-entry'', the enclave is re-entered every single time an application calls the cryptographic library on the same core as that of the application.
We see very little difference (2.9\%) between the first two setups, and a more important 27.6\% slowdown between the first and third scenario.
Indeed, our cryptographic algorithm is compute-heavy, making the communication overhead negligible in the second scenario.
Nevertheless, context-switching in and out of an enclave is expensive as it requires flushing the microarchitectural state, which makes the slowdown more important in the re-entry scenario.
These two secure solutions highlight an interesting tradeoff between resource occupation (one core and a DRAM region vs. only a DRAM region) and performance.

\noindent\textbf{Private Inference}
\label{sec:evalprivateinference}
We evaluate the following scenario where a remote user wants to do perform some private inference task on a neural network evaluated on an untrusted server.
We pace a simple neural network model inside of one of our enclaves.
We use onnx2c \cite{onnx2cgithub} to convert a ONNX file describing a handwritten digit recognition model (trained on MNIST) to standard C code and run private inference inside of our enclave.
The user uses our remote attestation mechanism to verify the enclave identity and establish a shared secret with it.
The user is then able to send 1024 encrypted requests to the enclave that are passed along to the enclave by the untrusted OS.
Upon reception, the enclave decrypts the ciphertext and performs the inference task on the private input.
It then sends back an encrypted output to the remote user.
In total, we need to implement 227 lines of straightforward C, which is to be compared to the 1536 LOC of the neural network and the 410 LOC of the AES library and the 4386 LOC of ED25519 we now need to include inside of the enclave for encryption and key exchange.
With the enclave code added, the TCB is now 21K LOC and the final enclave binary measures 54KB.

We now compare the performance between our setup and an insecure baseline where the model is directly evaluated outside of an enclave and where the LLC is not partitioned. We do consider that incoming traffic from the remote user is encrypted and needs to be decrypted first.
In both setups, runtime is split 55\% and 45\% between AES operations and inference.
We see negligible (0.76\%) overhead for running a model inside of the enclave, with all of it coming from the necessity to copy the ciphertext in enclave private memory before performing decryption.
This is not surprising as the rest of the computation happens in enclave private memory and should not suffer from any slowdown.
This doesn't account for the time required to set up the enclave, which is a one-time cost but represents 23.3\% of the cost of performing 1024 inference.

\subsection{Evaluating Secure Shared Memory}

We first demonstrate that naive public memory sharing is vulnerable to transient execution attacks by mounting an attack on an insecure version of Citadel where speculation on shared memory is allowed. 
The attack uses the LLC as a side channel and use a specter gadget located in enclave code to leak arbitrary enclave memory to untrusted code ruining on another core. Anecdotally, we were not able to successfully mount the attack on our secure prototype. 

\label{sec:evalsharedmemory}

\begin{table}[]
    \newcolumntype{C}[1]{>{\centering\let\newline\\\arraybackslash\hspace{0pt}}m{#1}}
     \begin{NiceTabular}{|C{1.4cm}|C{1.8cm}|C{2.3cm}|}[hvlines, corners=NW]
        \CodeBefore
        \rectanglecolor[gray]{0.9}{1-1}{1-3}
        \rectanglecolor[gray]{0.9}{1-5}{2-5}
        \Body
            \multicolumn{1}{c|}{} & Signature & Hash \\
            \Block{}{Citadel} & 1.027 & 1.464 \\
            \Block{}{- Burst} & 1.194 & 4.973 \\
            \Block{}{++ Burst} & 1.654 & 4.959 \\
    \end{NiceTabular}
    \vspace{0.40cm}
    \caption{Normalized runtime for signing and hashing three thousand random messages sent by the OS to an enclave. Runtime is normalized to the insecure baseline where the library is directly linked to the testbench program. We compare Citadel with two other configurations where the library code is not annotated for burst accesses to shared memory (- Burst) or all annotated (++ Burst).}
    \label{tab:resultcrypto}
\end{table}

We then evaluate the perfromance overhead of our comunication mechanism.
We use a similar setup as Section \ref{sec:evalcrypto} but only measure time spent computing the cryptographic function.
The enclave performs $3000$ cryptographic operations (signature or cryptographic hash) over random inputs provided by the OS.
We compare three setups.
First, our final design with burst mode enabled around shared memory burst sequences (Citadel).
Then Citadel without annotating code for shared memory burst sequences (- Burst).
Finally, another (insecure) configuration where we annotate all the code to be run in burst mode (++ Burst).
Our results are presented in Table \ref{tab:resultcrypto}.
We report the runtime normalized to the insecure baseline (directly linked library and no LLC partitioning).
We observe an overhead of 2.7\% on the signature benchmark and 46.6\% for the hashing benchmark.
This is because the hashing application accesses shared memory much more frequently compared to the signing application that is more computationally bounded.
The hashing function is more affected by the loss of the L1 cache.
This difference between the benchmarks is also highlighted by our other results, where not annotating the code for burst shared accesses has an  $397.3\%$ overhead for the hash application and a $19.4\%$ for the signing applicaiton.
These results also demonstrate the efficiency of our burst mechanism when used properly.

We also investigate the performance of using this mode on the whole enclave code instead of the limited code segment when the enclave communicates.
On our benchmarks, we observe that the performance overhead goes up to $65.4\%$ for the signing benchmark and $395.9\%$ for the hashing benchmark. 
We investigate the origin of this overhead and realize that miss speculation due to the deactivation of the RAS generates a lot of random load that pollutes the TLB, triggers many page walks and terribly slows down the virtual memory subsystem.
This is due to many small functions being linked one after the other in the binary.
Some of these functions start with loads using generic registers as addresses (for instance when accessing a pointer passed as an argument to the function).
When encountering a \texttt{ret} instruction, the processor will first speculate that the following instructions should be executed.
In this speculative window, it will start running the instructions of the function linked right after and perform loads or stores to random addresses from registers containing garbage.
We confirm this hypothesis by instrumenting the binary and observing that when the RAS is reactivated (but not other predictors), the number of TLB misses, page walks and the overall execution time go down significantly.
As a result, we recommend a hybrid approach where the burst mode is only used when the enclave wants to communicate.

\subsection{Evaluating LLC Dynamic Partitioning}
\label{sec:evalllcpartitioning}

We evaluate the advantage of dynamically partitioning the LLC. To do so, we evaluate two benchmarks, our cryptographic enclave performing signing operations, and running 2000 iterations of the Coremark \cite{gal2012exploring} benchmark wrapped inside an enclave.
We compare three different LLC configurations:
an insecure baseline where the LLC is not partitioned, a secure configuration where the LLC is statically partitioned into uniform slices of 16KB (MI6 scheme) and a last configuration where we use our scheme to partition the LLC non-uniformly.
For the non-uniform partition, we split the 1MB of LLC as follows: we give 256KB to the enclave, 16KB to the SM, 128KB to the OS, 128KB to the region used for shared memory 128KB to the region storing page tables and 1KB each to every other region.
For the Coremark benchmark, the uniform partitioning induces a 3.2\% overhead, while the non-uniform partitioning induces 0\% overhead. 
This indicates that the memory footprint of Coremark is smaller than 16KB.
For the Signature benchmark, we observe a 67.8\% overhead for the uniform partitioning and 0.08\% overhead for the non-uniform partitioning.
These last results highlight the benefits of dynamic partitioning, which provides security for no cost, whereas static partitioning more than doubles the total runtime of our application.
This is especially important as the benefits of dynamic partitioning are felt by both the enclave and untrusted applications.

We also measure the time it takes to modify LLC partitioning.
We measure that the operation takes on average 1M cycles.
This is a very costly operation, but it only happens very infrequently and not while programs are running.
Among these, 98\% of the cycles are spent flushing the LLC.
As a result, changing the layout of less regions is much faster: we measure an average of 2.5K cycles to flush a single LLC set containing 16 ways. 
The bigger the number of LLC set indexes affected, the longer the operation will take.

\subsection{Evaluating the Relaxed Flushing Policy}
\label{sec:evalrealxedflush}

\begin{figure}[t]
\centering
\includegraphics[width=\columnwidth]{figures/timer_interrupt_coremark.png}
\caption{Normalized Execution Time for Coremark in the presence of periodic timer interrupts. MI6 flushed on every interrupt; Citadel does not.}
\label{fig:coremark_timer}
\end{figure}

To evaluate the performance of the relaxed flushing policy, we run Coremark while triggering periodic timer interrupts at different time intervals to model frequent traps.
We do not perform any computation inside the trap handler and only measure the overhead due to context switching.
Results are provided in Figure \ref{fig:coremark_timer}.
By testing a range of interrupt periods, we confirm that flushing on each trap like in MI6 results in significant overheads.
However, the cost is greatly reduced when using the relaxed flushing policy, with a maximum of 7\% overhead measured for interrupts every 0.4ms.